\hypertarget{appendix-45303d03796a}{}\label{referencestart:appendix-45303d03796a}
\section{Integral promotion}\label{appendix-45303d03796a-integral-promotion}

\hypertarget{appendix-45303d03796a-integral-promotion}{}

\textbf{Curriculum depth:} Developed · \textbf{First seen:} C001 · \textbf{Latest development:} C002

\begin{allmightroledefinition}

\subsection{Definition and mechanism}\label{appendix-45303d03796a-definition-and-mechanism}

\hypertarget{appendix-45303d03796a-definition-and-mechanism}{}

Narrow integral operands promote to \BookInlineCode{int} when it represents their entire range, otherwise to \BookInlineCode{unsigned\ int}; the operator acts at that promoted width.

\end{allmightroledefinition}

Narrow integer types become int if int represents every original value; otherwise they become unsigned int. Promotion happens before the operation, not only on assignment. On the usual host, uint8\_t\{250\}+uint8\_t\{10\} is int 260; converting afterward to uint8\_t yields 4.

\subsection{Prerequisites}\label{appendix-45303d03796a-prerequisites}

\hypertarget{appendix-45303d03796a-prerequisites}{}

\begin{itemize}
\tightlist
\item
  \hyperref[appendix-cc04044b5782]{Fixed-width integer types} --- When supplied, std::uint32\_t has exactly 32 bits; uint\_least32\_t has at least 32; uint\_fast32\_t is chosen for speed with at least 32.
\item
  \hyperref[appendix-f414bd4ed05c]{Interpretation contract} --- For \BookInlineCode{0xAD}, unsigned interpretation gives 173; eight-bit two's complement gives −83.
\end{itemize}

\subsection{Used by}\label{appendix-45303d03796a-used-by}

\hypertarget{appendix-45303d03796a-used-by}{}

\begin{itemize}
\tightlist
\item
  \hyperref[appendix-5e2d663adf7f]{Usual arithmetic conversions} --- uses integral promotion as part of its representation or reasoning.
\end{itemize}

\subsection{Recall and next depth}\label{appendix-45303d03796a-recall-and-next-depth}

\hypertarget{appendix-45303d03796a-recall-and-next-depth}{}

Explain integral promotion without the label, work the example above, then state the assumption whose removal would change the result.

\textbf{Pending:} enumeration promotions, overload resolution, and implementation-specific ranks.

\subsection{Source and navigation}\label{appendix-45303d03796a-source-and-navigation}

\hypertarget{appendix-45303d03796a-source-and-navigation}{}

\hyperref[01-foundations-i-ch001-from-a-signal-to-an-interpretation]{C001 teaching chapter} · \hyperref[01-foundations-i-ch001-arithmetic-becomes-a-policy]{C002 development} · \href{../01-Bits-Programs-and-Reliable-Systems-P03/main.pdf\#nameddest=appendix-fd090d982a15}{Part definitions and mechanisms (Bits, Programs, and Reliable Systems, Part 3, Reference notes: Foundations)} · \href{../07-Vision-XR-and-Playable-Worlds-P01/main.pdf\#nameddest=appendix-4ad47bdd269a}{Dependency Map (Vision, XR, and Playable Worlds, Part 1, Reference notes: Dependency Map)}
